% !PRINTED 61
\weberpage{61}

If we wished to assume $l_1(u)$, $l_2(u)$ to be constants, then $T(u)$ would be an integral doubly periodic function of the second kind, and consequently, by \S\,16, 6., an exponential function, from which no doubly periodic functions can be formed. The simplest assumption that then remains to us is that $l_1(u)$, $l_2(u)$ are \gesperrt{linear functions} of $u$, and this includes as a special case that one of the two functions $l_1$, $l_2$ is constant.

We therefore set up the following definition:

1. \gesperrt{A $T$-function, $T(u)$, is an integral function of $u$ which satisfies the following two equations:}

\weq{6}{\begin{aligned}
T(u+\omega_1) &= e^{-\pi i [a_1(2u+\omega_1)+b_1]}\,T(u),\\
T(u+\omega_2) &= e^{-\pi i [a_2(2u+\omega_2)+b_2]}\,T(u).
\end{aligned}}

Here $a_1$, $b_1$, $a_2$, $b_2$ are constants, that is, quantities independent of $u$. The periods $\omega_1$, $\omega_2$ are here taken throughout as constants given once and for all, whose ratio $\omega_2:\omega_1$ has a positive imaginary part.

The factors

\[
\begin{aligned}
e_1 &= e_1(u) = e^{-\pi i [a_1(2u+\omega_1)+b_1]},\\
e_2 &= e_2(u) = e^{-\pi i [a_2(2a+\omega_2)+b_2]}
\end{aligned}
\]
% [sic] German source prints 2a in the second exponent; mathematically expected 2u.

are called the two \gesperrt{periodicity factors} of the function $T$.

The periodicity factors do not change when $b_1$, $b_2$ are altered by \gesperrt{even integers}. That functions of this kind really exist will appear only in the further course of the investigation.

2. \gesperrt{If the function $T(u)$ vanishes at any point, then it vanishes also at all congruent points. If it vanishes at $m$ points of the period parallelogram, it is said to be of the $m$th order.}

From this definition it follows at once:

3. \gesperrt{By multiplication of two $T$-functions $T$, $T'$ of order $m$, $m'$ one obtains a $T$-function of order $m+m'$. The periodicity factors of the product $T$, $T'$ are the products of the periodicity factors of $T$ and $T'$.}
